\section{Wyniki}
Model z takimi ograniczeniami \textbf{nie znajduje} rozwiązania \\
Problemem jest tu ogranicznie mówiące o tym że należy dostaczyć co najmniej \textbf{5000} ton \textbf{każdego} produktu \\
To ograniczenie jest niewykonywalne gdyż nawet zakładając dostarczenie maksymalnej liczby ton surowców (12000 + 8000 = 20000) ton, i ignorując ograniczenie przepustowośći przygotowalni 16000 ton, nie jesteśmy w stanie wytworzyć wystarczająco półproduktu D1 aby wytworzyć wystarczająco wyrobu W1 \\ 
gdybyśmy całość dostępnych surowców (znowu, ignorując przepustowość przygotowalni), zużyli na wytwarzanie tylko i wyłącznie półproduktu D1 z którego potem uzyskujemy wyrób W1, otrzymalibyśmy: 
\[ (12000 + 8000) * 0.1 = 20000 * 0.1 = 2000 \]
maksymalna liczba ton wyrobu W1 którą możemy uzyskać, ignorując przepustowość przygotowywalni oraz wymagania dotyczące W2 wynosi 2000 ton, czyli mniej niż 5000 wymaganych ton \\
W związku z tym wyniki przedstawione niżej dotyczącą przypadku w którym \textbf{ignorujemy} ograniczenie związane z minimalną liczbą ton każdego produktu który musimy dostarczyć (ograniczenie 20) \\
\begin{figure}[H]
\caption{}
\begin{tabular}{|c|c|c|c|c|c|c|c|c|c|c|c|}
\hline
 & 	z & S1 & S2 & G1 & P1 & P2 & O2 & D1 & D2 & W1 & W2 \\
\hline
z &	0 & 12000 & 8000 & 0 & 0 & 0 & 0 & 0 & 0 & 0 & 0 \\
\hline
S1 &	0 & 0 & 0 & 12000 & 0 & 0 & 0 & 0 & 0 & 0 & 0 \\
\hline
S2 &	0 & 0 & 0 & 0 & 0 & 2000 & 6000 & 0 & 0 & 0 & 0 \\
\hline
G1 &	0 & 0 & 0 & 0 & 12000 & 0 & 0 & 0 & 0 & 0 & 0 \\
\hline
P1 &	0 & 0 & 0 & 0 & 0 & 0 & 0 & 1200 & 10800 & 0 & 0 \\
\hline
P2 &	0 & 0 & 0 & 0 & 0 & 0 & 0 & 200 & 1800 & 0 & 0 \\
\hline
O2 &	0 & 0 & 0 & 0 & 0 & 0 & 0 & 0 & 0 & 0 & 6000 \\
\hline
D1 &	0 & 0 & 0 & 0 & 0 & 0 & 0 & 0 & 0 & 1400 & 0 \\
\hline
D2 &	0 & 0 & 0 & 0 & 0 & 0 & 0 & 0 & 0 & 0 & 12600 \\
\hline
W1 &	0 & 0 & 0 & 0 & 0 & 0 & 0 & 0 & 0 & 0 & 0 \\
\hline
W2 &	0 & 0 & 0 & 0 & 0 & 0 & 0 & 0 & 0 & 0 & 0 \\
\hline

\end{tabular}
\end{figure}

\[ a_1 = 9613 \]
\[ a_2 = 5341 \]
\[ b_1 = 5910 \]
\[ b_2 = 3651 \]
\[ w_1 = 56 \]
\[ c_{P1} = 667 \]
\[ e_{P2} = 80 \]
\[ e_{O2} = 240 \]
\[ e_1 = 0 \]
\[ e_2 = 0 \]
\[ l = 94 \]
Ponownie, wyniki powyżej są prawdziwe wyłącznie dla przypadku w którym \textbf{ignorujemy ograniczenie numer 20}

